% Figura propia (diagrama o gráfico generado por figura.py).
% \includegraphics sin extensión: toma el .pdf si existe; si no, el .png.
\begin{figure}[htbp]
\caption{Arquitectura de la solución propuesta}
\label{fig:arquitectura}
\centering
\includegraphics[width=\linewidth]{figuras/out/arquitectura}
\nota{Elaboración propia.}
\end{figure}

% Figura tomada de otra fuente: la cita es obligatoria. Si la modificaste: "Adaptado de".
\begin{figure}[htbp]
\caption{Estructura de una red neuronal de propagación hacia atrás}
\label{fig:bp}
\centering
\includegraphics[width=0.8\linewidth]{figuras/out/estructura-bp}
\nota{Tomado de \textcite{li2016}.}
\end{figure}

% Captura del sistema (sección de implementación). \fbox da un borde fino.
\begin{figure}[htbp]
\caption{Panel web de atención de conversaciones}
\label{fig:panel}
\centering
\fbox{\includegraphics[width=0.95\linewidth]{figuras/out/panel-conversaciones}}
\nota{Captura del sistema implementado. Elaboración propia.}
\end{figure}

% En el texto, siempre antes de la figura:
%   "... como se observa en la Figura~\ref{fig:arquitectura}, ..."
